\documentclass[10pt]{article}
\usepackage[T1]{fontenc}
\usepackage[margin=0.9in]{geometry}
\usepackage{amsmath,amssymb,mathtools}
\usepackage{booktabs,array}
\usepackage{xcolor}
\usepackage[hidelinks]{hyperref}
\usepackage{microtype}
\usepackage{fancyhdr}
\definecolor{navy}{RGB}{24,55,95}
\pagestyle{fancy}
\setlength{\headheight}{14pt}
\fancyhf{}
\lhead{Finite polar groups}
\rhead{Point actions and exact verification}
\cfoot{\thepage}
\newcommand{\Sym}{\operatorname{Sym}}
\newcommand{\Fix}{\operatorname{Fix}}
\newcommand{\F}{\mathbb F}
\newcommand{\Sg}{\mathcal S}
\title{\color{navy}\textbf{Finite Polar-Space Permutation Actions}\\[2mm]
\large Unitary, symplectic, and orthogonal $\sigma$-MGFs}
\author{Proof and computational verification companion}
\date{17 July 2026}

\begin{document}
\maketitle

\begin{abstract}
This note specializes the exact permutation $\sigma$-MGF to isotropic and
singular point actions of finite unitary, symplectic, and orthogonal groups.
Witt extension proves the stable three-orbit classification of pairs.  The
companion marimo notebook constructs $Sp_4(2)$, $O_4^+(2)$, and $U_4(2)$ as
finite-field matrix groups, induces their point actions, and independently
compares direct orbit counts, cycle products, fixed-point reconstruction, and
matrix determinants.  Every tested coefficient agrees exactly.
\end{abstract}

\section{Exact formula for a polar action}
Let a finite classical group $G$ act on a finite set $S$ of isotropic or
singular one-spaces.  For $g\in G$, let $c_d(g;S)$ denote the number of
$d$-cycles of its action on $S$.  The permutation-matrix eigenvalues on a
$d$-cycle are the $d$th roots of unity, so
\begin{equation}
 \boxed{\Sg_{G,S}(t_\bullet)=\frac1{|G|}\sum_{g\in G}
 \prod_{i,d\geq1}(1-t_i^d)^{-c_d(g;S)}.}
\end{equation}
Moreover,
\begin{equation}
 \Fix_S(g^r)=\sum_{d\mid r}d\,c_d(g;S),\qquad
 c_d(g;S)=\frac1d\sum_{r\mid d}\mu(d/r)\Fix_S(g^r).
\end{equation}
Thus one may compute (1) purely from the number of isotropic or singular
eigenlines of every power $g^r$.  Equivalently, for a parabolic point
stabilizer $P$, this fixed-point count is
$(\operatorname{Ind}_P^G\mathbf1)(g^r)$.

For $\tau=(\tau_1,\ldots,\tau_s)$, expanding the determinants gives
\begin{equation}
 [m_\tau]\Sg_{G,S}=\dim\left(
 \bigotimes_{b=1}^s\Sym^{\tau_b}\mathbb C[S]\right)^G.
\end{equation}
The tensor product has a basis consisting of tuples of multisets of points.
Because $G$ permutes this basis, (3) is exactly the number of $G$-orbits on
such tuples.  This orbit count is the independent direct computation used
below.

\section{Why the degree-two coefficient is three}
For two isotropic or singular points in a nondegenerate polar space of
sufficient Witt rank, there are three possible geometric types:
\begin{enumerate}
 \item the two points coincide;
 \item they are distinct and perpendicular;
 \item they are distinct and nonperpendicular.
\end{enumerate}
Each relation is preserved by every isometry.  Conversely, an isometry
between either pair of generated subspaces extends to the ambient space by
Witt's extension theorem.  Hence each type is one orbit.  This applies to
both ordered pairs and unordered 2-multisets, proving
\begin{equation}
 \boxed{[m_{(2)}]\Sg=[m_{(1,1)}]\Sg=3.}
\end{equation}
It follows that the stable expansion begins
\[
 \Sg=1+m_{(1)}+3m_{(2)}+3m_{(1,1)}+O(\deg3).
\]
The extra perpendicularity relation already distinguishes these limits from
the Poisson/symmetric-group series in degree two.

\paragraph{Proposition (coefficientwise stability within a fixed polar tower).}
Fix the field, one classical type, and a total degree $D$.  For each tuple of
multisets of singular or isotropic points of total degree $D$, retain the
labelled span together with the restricted form and the multiplicities of the
labelled one-spaces; in characteristic-two orthogonal cases retain the
restricted quadratic form rather than only its polar bilinear form.  Two
configurations in the same ambient nondegenerate polar space are conjugate
precisely when these labelled formed spaces are isometric, because an
isometry of their spans extends by Witt's lemma.

Only finitely many labelled formed spaces of dimension at most $D$ exist over
the fixed finite field.  Let $m_0(D)$ be the maximum, over the types that occur
somewhere in the chosen tower, of the least Witt rank in which the type
embeds.  Every coefficient of total degree $D$ is therefore constant for Witt
rank at least $m_0(D)$.  Orthogonal plus, minus, and odd-dimensional towers
are treated separately.  This proves existence of a stable threshold but
does not claim that the displayed bound is sharp.  The degree-two value three
was proved directly above and does not depend on this proposition.

\section{The three exact matrix constructions}
All arithmetic in the program is exact integer arithmetic.

\paragraph{$Sp_4(2)$.}
We enumerate $GL_4(2)$ and retain exactly the matrices preserving
\[
 B(x,y)=x_1y_2+x_2y_1+x_3y_4+x_4y_3.
\]
The resulting group has order $720$.  In characteristic two every one-space
is isotropic, giving a permutation action on the 15 nonzero vectors (equally,
the 15 projective points).

\paragraph{$O_4^+(2)$.}
We retain the matrices preserving the hyperbolic quadratic form
\[
 Q(x)=x_1x_2+x_3x_4.
\]
The resulting group has order $72$ and acts on its 9 nonzero singular vectors.

\paragraph{$U_4(2)$.}
We implement $\F_4=\F_2[\alpha]/(\alpha^2+\alpha+1)$ and the Hermitian form
\[
 h(x,y)=\sum_{j=1}^4\overline{x_j}y_j,\qquad \overline z=z^2.
\]
Every ordered orthonormal basis is enumerated recursively, producing exactly
\[
 |U_4(2)|=2^6\prod_{i=1}^4(2^i-(-1)^i)=77{,}760
\]
matrices.  Normalizing each nonzero vector by its first nonzero coordinate
gives 45 isotropic projective points.  The scalar center has order three, so
the faithful induced permutation image has order $25{,}920$.

\section{Verification strategy and results}
For each action, the notebook:
\begin{enumerate}
 \item constructs the full finite-field matrix group and induced permutation;
 \item directly enumerates orbits on ordered point pairs and 2-multisets;
 \item extracts the same coefficients from the cycle product in (1);
 \item reconstructs cycle counts from fixed points of powers using (2); and
 \item numerically compares direct permutation-matrix determinants with the
 cycle products at $(t_1,t_2)=(0.07,0.11)$.
\end{enumerate}
The regression module also compares a degree-three labelled-formed-space
orbit count in two consecutive ranks of each implemented tower.  This tests
the restricted-form encoding; it is not used as a proof of the proposition.

\begin{center}
\begin{tabular}{@{}lcrrr@{}}
\toprule
Action & $\tau$ & direct orbits & cycle formula & claimed\\
\midrule
$Sp_4(2)$ on 15 isotropic points & $(2)$ &3&3&3\\
 & $(1,1)$ &3&3&3\\
$O_4^+(2)$ on 9 singular points & $(2)$ &3&3&3\\
 & $(1,1)$ &3&3&3\\
$U_4(2)$ on 45 isotropic points & $(2)$ &3&3&3\\
 & $(1,1)$ &3&3&3\\
\bottomrule
\end{tabular}
\end{center}

Every equality in the table is exact.  All cycle counts in the tested
$Sp_4(2)$ and projective $U_4(2)$ actions are also recovered exactly from
fixed points of powers.  For the symplectic action, the direct determinant
average and cycle-product average both equal $1.304867127368$, with zero
difference at the displayed floating-point precision.

\section{Representations requiring additional choices}
The Weil-representation paragraph is governed by the general identity
\[
 \det(I-t\rho(g))^{-1}=\exp\left(
 \sum_{r\geq1}\frac{\chi_\rho(g^r)}{r}t^r\right),
\]
but it does not specify one numerical representation.  A reproducible matrix
test requires a choice of $q$ and $n$, an additive character, a linearization,
and whether $\rho$ is the full Weil representation or a constituent.  Even
characteristic requires a different construction.  Accordingly, no arbitrary
choice was silently inserted into the verification.

Likewise, the exceptional/Suzuki/Ree formula is an exact template for a
specified pair $(G(q),P(q))$, not a single representation.  A numerical test
requires the concrete group, field size, and parabolic action.  Once supplied,
the same fixed-point, cycle, determinant, and orbit pipeline applies without
change.

\paragraph{Conclusion.}
The three classical polar families realize exactly the predicted three pair
types in small but stable-rank examples.  The agreement of full matrix-group
enumeration, direct configuration orbits, cycle extraction, M\"obius
reconstruction, and determinants validates both the formula and the proposed
verification method.

\end{document}
